\documentclass[10pt]{article}%
\usepackage[hidelinks]{hyperref}
\hypersetup{
  pdfauthor={Joshua N Grant},
  pdftitle={CV for Joshua N. Grant}
}
\author{Joshua N Grant}
\title{CV for Joshua N. Grant}
\usepackage{etoolbox}
\patchcmd{\thebibliography}{\section*{\refname}}{}{}{}
\usepackage{tikz}
\usepackage{xcolor}
\usepackage[left=1cm,top=1cm,right=1cm,bottom=1cm,nohead,nofoot]{geometry}
\usepackage{tabularx}
\usepackage{makecell}
%\usepackage{palatino}
\usepackage{newpxtext,newpxmath}
\usepackage{pifont}
\usepackage{enumitem}
\usepackage{sectsty}
\usepackage{fontawesome}
\bibliographystyle{plain}
\definecolor{fillheader}{HTML}{FFFFFF}
\definecolor{header}{HTML}{000000}
\definecolor{headercolor}{HTML}{e3e3e3}
\definecolor{linkcolor}{HTML}{3e3e3e}
\usetikzlibrary{positioning,shapes,shadows,arrows,backgrounds,calc,shadows.blur}
\newcolumntype{M}{>{\raggedright\arraybackslash}X}
% custom macros for the resume
\newcommand{\name}[2]{\makecell[l]{{\Large #1 \huge \textbf{#2}}}}
\newcommand{\objective}[1]{
    \begin{minipage}[ht]{0.2\linewidth}
    \section*{\faBullseye{} Objective}
    \end{minipage}
    \begin{minipage}[ht]{0.8\linewidth}
        \centering #1
    \end{minipage}
}
\newcommand{\contact}[6]{\section*{\faUser{} Contact} \vspace{-0.5em} {\begin{minipage}[ht]{.33\linewidth}%
    \faMapMarker \, #1 \\ \faPhone \, #2 \end{minipage}%
    \begin{minipage}[ht]{.66\linewidth}%
    \faEnvelope \, \href{mailto:#3}{#3} \\
    \faBold \, \href{https://#4}{#4} \\
    \faGithub \, \href{#5}{#5} \\ \faLinkedin \, \href{#6}{#6}
    \end{minipage}}
}
\newcommand{\skills}[3]{\makecell[l]{#1 \textbf{#2:}\hspace{0.22em} #3}}
% Macro for Projects Section
\newcommand{\projects}[3]{%
    \textbf{#1} \textemdash{} \textit{#2} \\ % Summary
    \textbf{Tech:} #3 % Technologies Used
    \vspace{-1em}
    \begin{center}
        \rule{0.7\textwidth}{0.2pt}
    \end{center}
}

\newcommand{\job}[4]{
    \makecell[l]{
        \faBriefcase{} \textbf{#1} \\
        \textit{#2} \\ #3 \\
        \begin{itemize}[topsep=-12pt,parsep=0em] \setlength\itemsep{0em} #4 \end{itemize} \\}
}
\newcommand{\education}[4]{\makecell[l]{\faGraduationCap{} \textbf{#1} \\ \textit{#2} \\ \textbf{#3}, #4 \\}}
\newcommand{\references}[1]{\section*{\faUserPlus{} References} \vspace{-0.5em} \makecell[l]{#1}}
\newcommand{\header}[1]{\section*{\faCircleInfo{} #1} \vspace{-0.5em}}
\newcommand{\headerline}{\noindent\rule{\textwidth}{0.4pt}}
% begin the document
\begin{document}
    \pagenumbering{gobble}
    \begin{minipage}[ht]{.33\linewidth}%
    % name section
    \name{Joshua N.}{Grant}
    \vspace{-0.5em}
    \makecell[l]{\large \textit{Geospatial Systems Architect}}
    \end{minipage}%
    \begin{minipage}[ht]{.66\linewidth}%
    % contact section
    \contact{12918 Sanderling Lane\\ Knoxville, TN 37922 \\ United States of America}{+1 (865) 803 3495}
        {jngrant@live.com}
        {notjustadatum.blogspot.com}
        {github.com/sempervent}
        {linkedin.com/in/joshuanagrant}
    \end{minipage}%
    \par\noindent\rule{\textwidth}{0.4pt}
    % Objective Section
    \objective{To take theories, ideas, research, and experiments and turn them into production-scalable solutions \\ while developing and supporting applications and data workflows on cloud infrastructure.}
    \headerline
    % Skills Section
\section*{\faCogs{} Skills}
    \begingroup
    \setlength{\extrarowheight}{-2pt}
    \renewcommand{\arraystretch}{0.5}
    \begin{tabularx}{\linewidth}{@{}>{\raggedright\arraybackslash}p{0.2\linewidth}X@{}}
        \vspace{-1.5em}
        \skills{\faCode}{Programming Languages}{Python, R, JavaScript, SQL, Bash, PHP, Rust \LaTeX}
        \vspace{-1.5em}
        \skills{\faCloud}{Cloud Platforms}{AWS (S3, Lambda, EC2, RDS, Fargate), GCP, Azure}
        \vspace{-1em}
        \skills{\faDatabase}{Databases}{PostgreSQL, MongoDB, MySQL, SpatiaLite, Snowflake, SQLite, DuckDB, Redis, RabbitMQ, MQTT}
        \vspace{-1em}
        \skills{\faGears}{Data Engineering}{Airflow, Prefect, Kafka, Scala, Hadoop, TimescaleDB, NiFi, ETL Pipelines}
        \vspace{-1em}
        \skills{\faMapMarker}{Geospatial Tools}{ArcGIS, GeoPandas, Geoparquet, QGIS, GDAL, PostGIS, Martin, Pyproj, Rasterio, MapBox, Leaflet}
        \vspace{-1em}
        \skills{\faPlay}{DevOps}{Docker, Kubernetes, Ansible, GitHub Actions, AWS CodePipeline}
        \vspace{-1em}
        \skills{\faDashboard}{Data Visualization}{Shiny, Leaflet, Kibana, matplotlib, Plotly, Grafana}
        \vspace{-1em}
        \skills{\faFlask}{Frameworks and APIs}{FastAPI, Flask, Django, Swagger, REST API Development, Electron}
        \vspace{-1em}
        \skills{\faBolt}{Machine Learning}{TensorFlow, Scikit-learn, Data Preprocessing, Model Deployment}
        \vspace{-1em}
        \skills{\faGit}{Version Control}{Git (GitHub, GitLab, Bitbucket)}
        \vspace{-1em}
        \skills{\faUsers}{Soft Skills}{Team Leadership, Agile Methodologies, Cross-Functional Collaboration, Written \& Verbal Communication, \\ Problem Solving}
        \vspace{1em}
    \end{tabularx}
    \endgroup
    \headerline
    % Projects Section
    \section*{\faPuzzlePiece{} Projects}
    \projects{Final Fantasy Football}{A semantic learning algorithm for fantasy football}{Python, Scikit-learn, Pandas, NumPy, Matplotlib, Requests, BeautifulSoup, Flask, Docker, Docker Compose, Airflow}
    \projects{Where I've Been}{A web application to visualize travel history}{Python, Flask, Docker, Docker Compose, PostgreSQL, Leaflet, Vue, Bootstrap}
    \projects{This Is A Casino}{A web application to visualize and trade stocks via semantic learning, reinforcement learning, and deep learning}{Python, Flask, Docker, Docker Compose, PostgreSQL, Leaflet, React, Bootstrap}
    \projects{Genesis}{A python audio input game to create a universe from a sound}{Python, pygame}
    \projects{Decentralized Content Reward System}{A decentralized content reward system for the web}{Python, Rust, Flask, Docker, Docker Compose, PostgreSQL, Leaflet, React, Bootstrap}
    \projects{Cosmic Architect}{A pygame game where you compete to build the best planet}{Python, pygame}
    \projects{maw}{A cli for quickly combining tabular data into parquet}{Rust}
    \headerline
    \section*{\faCalendar{} Experience}
    \begin{tabularx}{\linewidth}{M}%
        \makecell[l]{\textbf{Geospatial Systems Architect}} \\
        Oak Ridge National Laboratory, Oak Ridge, TN \textit{2023\textemdash{}Present} \\
        \begin{itemize}[topsep=-12pt,parsep=0em]
            \setlength\itemsep{0em}
            \item Architecting and maintaining a geoparquet-based data warehouse for storage and retrieval of risk components %
            \item Designing, architecting, and maintaining simultaneous deployment of a Geospatial Decision Support Tool for both web and desktop environments %
            \item Architecting a geospatial data pipeline, using Airflow for the ingestion, storage, and analysis of raster data in a cloud environment %
            \item Developing geospatial applications for the tracking and monitoring of high-sensitivity parcels using IoT technologies, Kafka, TimescaleDB, and Airflow running in AWS %
            \item Designing and architecting a geospatial data pipeline, database, frontend, and backend for the ingestion, storage, and retrieval of clock drift timeseries data using Airflow %
            \item Developing and maintaining a geospatial data warehouse and pipeline for the ingestion, storage, and retrieval of raster data with data ingestion controlled by Prefect %
            \item Architecting a geospatial database capable of handling large-scale raster data and aggregating zonal stats to a variety of vector geometries %
        \end{itemize} \\
        \makecell[l]{\textbf{Senior Software Engineer \dash Data Engineering}} \\
        Bold Penguin, Columbus, OH (remote) 2022\textemdash{}2023 \\
        \begin{itemize}[topsep=-12pt,parsep=0em]
            \setlength\itemsep{0em}
            \item Designed a more centralized system for development and deployment using GitHub Actions, AWS, and Python modules %
            \item Supervised a team effort to migrate deployments, flows, and tasks from Prefect 1 to Prefect 2 %
            \item Developed Prefect tasks to run on either EC2 or Fargate instances based on conditions such as file size %
            \item Created EC2 and Fargate microservices for assistance in document conversion and data extraction %
            \item Designed and maintained REST APIs to be up-to-date with stakeholder and client demands %
            \item Designed and implemented a scalable message queue (MQ) system to handle high load for a data extraction service, resulting in improved performance and increased efficiency of data processing operations %
            \item Configured GitHub Actions and AWS CodePipeline to automate the build, test, and deployment process of the projects, ensuring consistent and reliable delivery of software updates %
            \item Created and oversaw the deployment of multiple Python and JavaScript projects using AWS CodePipeline, resulting in streamlined, automated, and efficient project delivery %
        \end{itemize} \\
        \makecell[l]{\textbf{Data Engineer}} \\
        Oak Ridge National Laboratory, Oak Ridge, TN \textit{2019\textemdash{}2022} \\
        \begin{itemize}[topsep=-12pt,parsep=0em]
            \setlength\itemsep{0em}
            \item Served as support staff for a Python-based web scraping project, implementing an ETL workflow from a brittle-based cron system to a more robust Airflow system deployed in a Kubernetes cluster and in the cloud %
            \item Developed and maintained full stack docker images using Python, R, Airflow, and Shiny to automate the acquisition of spatial-temporal data for the WSTAMP project for storage in an Hadoop Data Warehouse %
            \item Led the development team for a COVID-19 tracking project using both R and Python and downstream model creation and analysis %
            \item Developed containerized Airflow ETL workflows, Python microservices, and REST APIs using Swagger for data exchange with PostgreSQL, AWS S3 \& RDS, and MongoDB %
            \item Automated web-scraping and social media ingestion via containerized model creation for misinformation classification for implementation in advanced statistical models and network analysis and ingestion into ElasticSearch and visualization in Kibana %
            \item Served as Software Quality Assurance board chair for the Geospatial Science and Human Security Division %
            \item Participated as a member of the ORNL Data Council, which discusses best practices in lab-wide data management and architecture for sharing of lab-generated data %
        \end{itemize} \\
        \makecell[l]{\textbf{Data Scientist II Subcontractor}} \\
        Oak Ridge National Laboratory, Oak Ridge, TN \textit{2018\textemdash{}2019} \\
        \begin{itemize}[topsep=-12pt,parsep=0em]
            \setlength\itemsep{0em}
            \item Automated data ingestion of the WSTAMP Web Application via the use of R, Python, Airflow, and Docker %
            \item Authored an R package to unify geography between online data APIs %
            \item Developed a SpatiaLite SQLite database to ease lookup of geographic data %
            \item Developed a zonal statistics application in Python that takes a raster image and shapefile and computes the statistics over the shapes stored in the shapefile %
        \end{itemize} \\
    \end{tabularx}
    \begin{tabularx}{\linewidth}{M}%
        \makecell[l]{\textbf{Owner \& Contractor-For-Hire}} \\
        Specrabella, Knoxville, TN \textit{2018\textemdash{}2018} \\
        \begin{itemize}[topsep=-12pt,parsep=0em]
            \setlength\itemsep{0em}
            \item Designed, developed, and deployed an interactive energy informatics leaflet-based web application integrated with Snowflake that provides real-time calculation from machine learning derived models %
            \item Consulted with clients on ground-up web application architecture, database and visualization needs%
            \item Developed data ingestion methods using cURL, PHP, and R to scrape energy data from the web %
            \item Streamlined training and development procedures using ansible, vagrant, PHP, and BASH %
            \item Designed and implemented a data pipeline using R, Python, and Snowflake to ingest, clean, and visualize energy data %
            \item Developed a web application to visualize energy data using R, Shiny, and Leaflet %
            \item Architected, designed, and maintained PostgreSQL databases for energy and financial data %
        \end{itemize}
        \end{tabularx}
    \begin{tabularx}{\linewidth}{M}%
        \makecell[l]{\textbf{Bioinformatician, Project Manager, \& Technical Sales Representative}} \\
        Microbial Insights, Knoxville, TN \textit{2016\textemdash{}2018} \\
        \begin{itemize}[topsep=-12pt,parsep=0em]
            \setlength\itemsep{0em}
            \item Established an automated pipeline using R, MySQL, Python, and Bash for NGS analysis, along with intranet website control written in JavaScript and Shiny  %
            \item Automated client-bound statistical calculations such as linear models, ANOVA, SOMs, clustering, PCA, and PCoA %
            \item Developed an ETL customer data visualization tool using R, PHP, JQuery, and MySQL to view qPCR results in the context of other samples' and dynamically selectable parameters %
            \item Optimized data delivery to clients via a custom R package and local shiny applications to quickly generate \LaTeX{} PDF reports  %
            \item Constructed, populated, and maintained an intranet wiki using PostgreSQL and PHP to aid in project management and customer service
        \end{itemize} \\
        \makecell[l]{\textbf{Graduate Research Assistant, Dr. Neal Stewart's Plant Biotechnology Lab}} \\
        University of Tennessee, Knoxville, TN \textit{2014\textemdash{}2016} \\
        \begin{itemize}[topsep=-12pt,parsep=0em]
            \setlength\itemsep{0em}
            \item Summarized statistical findings of cell suspensions using linear models, ANOVA, and PCA in Python and R %
            \item Evaluated suspension cultures via chemical and spectral processes for lignin formation and statistically analyzed and summarized my findings for inclusion in DOE reports %
            \item Collaborated on a novel single cell suspension and cryopreservation robotic system %
            \item Advanced and executed monocot genetic modifications for \textit{Panicum} spp., \textit{Oryza} spp., \textit{Sorghum} spp., \textit{Saccharum} spp., \textit{Zea mays}
        \end{itemize} \\
        \makecell[l]{\textbf{Laboratory Assistant, Dr. Neal Stewart's Plant Biotechnology Lab}} \\
        University of Tennessee, Knoxville, TN \textit{2012\textemdash{}2014} \\
        \begin{itemize}[topsep=-12pt,parsep=0em]
            \setlength\itemsep{0em}
            \item Developed automated statistical methodology for screening of lignin content and imaging of cell characteristics both \textit{in vivo} and \textit{in vitro} using Python%
            \item Extracted genomes from NCBI, cleaned and normalized the data, and performed exploratory data analysis using Python
            \item Complied with USDA-APHIS regulations regarding transgenic plants
            \item Implemented an \textit{E. coli} bioreactor for production of proteins
        \end{itemize} \\
        \makecell[l]{\textbf{Laboratory Assistant, Dr. Paris Lambdin's Biosystematics and Biological Control Lab}} \\
        University of Tennessee, Knoxville, TN \textit{2002\textemdash{}2012} \\
        \begin{itemize}[topsep=-12pt,parsep=0em]
            \setlength\itemsep{0em}
            \item Designed instructional modules for undergraduate and graduate level courses using HTML, CSS, JavaScript, and Flash
            \item Collected samples and data from field locations
            \item Assisted in community outreach programs such as Bloomsdays, Buggy Buffet, and 4-H Camps
            \item Performed forest coverage analysis using ArcGIS
        \end{itemize}
    \end{tabularx}
    \headerline
    % Education Section
\section*{\faUniversity{} Education}
\begin{tabularx}{\linewidth}{M}
    \education{Master of Science in Plant Sciences \textemdash Plant Molecular Genetics}{University of Tennessee}{Spring 2017}{GPA: 3.72/4.0}
    \vspace{0.5em}
    \education{Bachelor of Science in Plant Sciences \textemdash Biotechnology}{University of Tennessee}{Spring 2014}{GPA: 3.74/4.0 \textemdash Magna Cum Laude}
\end{tabularx}
    \headerline
    \section*{\faBook{} Publications}
    \nocite{*}
    \bibliography{bibliography}
    \par\noindent\rule{\textwidth}{0.4pt}
    \section*{\faUserPlus{} References}
    References available upon request.
\end{document}